\chapter{Program Architecture and Flow}
\label{ch:architecture}

\section{Module structure}

\begin{table}[H]
\centering
\caption{Core implementation modules.}
\label{tab:modules}
\begin{tabularx}{\textwidth}{@{}p{0.30\textwidth}X@{}}
\toprule
Module & Responsibility\\
\midrule
\texttt{picnn.py} & Partial ICNN architecture and reference normalization\\
\texttt{constitutive.py} & AD stress and full tangent extraction\\
\texttt{ngsolve\_umat.py} & Phase~1 integration-point bridge and Newton solve\\
\texttt{rve.py} & Periodic geometry and rigid kinematic utilities\\
\texttt{mixed\_matrix.py} & Taylor--Hood mixed matrix formulation\\
\texttt{rve\_solver.py} & Coupled periodic nonlinear RVE solution\\
\texttt{free\_particle.py} & Free finite rotation and rotational condensation\\
\texttt{homogenization.py} & Energy, stress, and Hill--Mandel work\\
\texttt{rve\_tangent.py} & Condensed macroscopic tangent\\
\texttt{rve\_dataset.py} & Deterministic sampling, checkpoints, and serialization\\
\bottomrule
\end{tabularx}
\end{table}

\section{End-to-end flow}

The implementation proceeds in the same order as the two project stages:
\begin{enumerate}
  \item \textbf{Phase 1 proof of concept:}
    \begin{enumerate}[label=\alph*)]
      \item sample analytical neo-Hookean labels and train the pICNN;
      \item use PyTorch automatic differentiation to obtain energy, stress,
            and the full tangent;
      \item transfer material values at NGSolve integration points and solve
            the displacement problem;
      \item pass the reference, response, tangent, equilibrium, patch, and
            matched-FE verification gates.
    \end{enumerate}
  \item \textbf{Phase 2 multiscale extension:}
    \begin{enumerate}[label=\alph*)]
      \item construct periodic RVEs and solve the mixed matrix and rigid
            particle equilibrium problem;
      \item homogenize energy and work-conjugate stress, and condense the
            consistent tangent;
      \item checkpoint the grouped dataset, train and evaluate the
            volume-fraction-conditioned pICNN, and run macroscopic NGSolve
            problems.
    \end{enumerate}
\end{enumerate}

\begin{figure}[H]
\centering
\small
\begin{tabular}{c}
\fbox{\parbox{0.88\textwidth}{\centering
  \textbf{Phase 1: prove the PyTorch--NGSolve material interface}\\
  Analytical compressible neo-Hookean data \(\rightarrow\) train the pICNN
  \(\rightarrow\) differentiate energy in PyTorch to obtain \(\Pstress\) and
  the full \(9\times9\) tangent.}}
\\[-1mm] \(\Downarrow\) \\[-1mm]
\fbox{\parbox{0.88\textwidth}{\centering
  Transfer integration-point material values through
  \texttt{IntegrationRuleSpace}; solve an NGSolve displacement problem with
  the consistent Newton tangent.}}
\\[-1mm] \(\Downarrow\) \\[-1mm]
\fbox{\parbox{0.88\textwidth}{\centering
  \textbf{Phase 1 gate:} reference state, analytical response, tangent,
  nonlinear solve, affine patch, and matched FE checks. Passing this gate
  precedes RVE data generation.}}
\\[1mm] \(\Downarrow\) \\[-1mm]
\fbox{\parbox{0.88\textwidth}{\centering
  \textbf{Phase 2: solve periodic rigid-particle RVEs}\\
  Mixed displacement--pressure PDMS matrix; exact rigid-sphere kinematics;
  equilibrium in displacement, pressure, and free rotation.}}
\\[-1mm] \(\Downarrow\) \\[-1mm]
\fbox{\parbox{0.88\textwidth}{\centering
  Homogenize energy and work-conjugate stress; compute the consistent
  macroscopic tangent by energy-Hessian condensation; checkpoint grouped
  data over deformation states and particle volume fractions.}}
\\[-1mm] \(\Downarrow\) \\[-1mm]
\fbox{\parbox{0.88\textwidth}{\centering
  Train and evaluate the volume-fraction-conditioned pICNN; apply boundary
  conditions in the macroscopic NGSolve problem and evaluate the learned
  constitutive response at its integration points.}}
\end{tabular}
\caption{Implemented two-stage workflow. The Phase~1 constitutive bridge was
verified before the Phase~2 RVE data-generation and homogenization workflow.}
\label{fig:flowchart}
\end{figure}
